\section{\ppm{[Ignore everything from here!]}Towards a proof of \ournotmain\ conjecture \ref{conj:onm0}}

The conjecture is proved in case $l=-1$. One proof uses (the quantum version of) Young's seminormal form, and hence hook lengths, 
so one strategy in general is to try to generalise this - giving more `geometric' 
bases for standard modules. 

\subsection{Formulation of a %\sout{the} 
Walk-Diagram Correspondence}
Here we aim to describe the emerging correspondence between walks in the Rollet graph and half-diagrams, describing a basis for a certain $J_{l,n}$ standard module. We will do this in baby steps, starting with the simple 1-cup, 1-D specht case:

\subsubsection{1-cup, 1-D Specht Case}
Here we will treat the case of the $J:=J_{l,n+2}$-module $\Delta:=\Delta_{n,(n)}$ (take $l+1\leq n$ WLOG), knowing that the case $\Delta:=\Delta_{n,(1^n)}$ is obtained by means of symmetry. 


As noted in (\ref{de:Rollet}):
Vertices in the Rollet graph  $\Roll_l$ are labelled by a pair $(\l, p)$, where $\l\vdash \min(p,l+2)$ for $p\in \N$, and there is an edge $(\l, p)-(\l', p')$ if and only if (up to interchanging the pairs) $p'=p+1$, and either $\l'$ can be obtained from $\l$ by adding a box, or $\l'=\l$.


\mdef Claim.  
We claim that there exists a correspondence between walks of length $n+2$ between $(\emptyset,0)$ and $((1^n),n)$ (in the Rollet graph of height $l$), and a basis of half diagrams for $\Delta$. Observe that for $\Delta$ a basis is simply described by the position of a single cup, since the corresponding specht module is 1D. We write $e_{i,j}$ for the basis element with a cup joining position $i$ and $j$. 

A walk is then described by a sequence of vertices $(\l_k, p_k)$, $k=-1, \dots, n,$ with $(\l_{-1}, p_{-1})=(\emptyset,0)$, $(\l_0, p_0)=((1),1), (\l_{n}, p_{n})=((1^n), n)$. It seems good to divide the walks into two types: 

A walk $(\l_k, p_k)$ is \textbf{ideal} if for all $k$, $\l_k \subset (1^n)$. We claim that \textbf{ideal} walks correspond to the height -1 diagrams as follows: Such walks are determined by the sequence $(p_l)$ and contain at least one repeated integer. Let $j$ be the maximal repeated integer which may take values $j=1,\dots, n$. The corresponding half diagram is $e_{j,j+1}$.

Now let $(\l_k,p_k)$ be a non-ideal walk. Since the sequence $p_k$ decreases exactly once, it follows that any $\l_k\not\subset (1^n)$, must be of the form $(l,1)$ for some $l$, and furthermore such a $\l_k$ is a partition of $p_k$ (i.e. $p_k$ does not exceed $l+2$). Therefore, a non-ideal sequence is uniquely determined by the following two integers:
\[i=\min\{k \ | \ \l_k\not\subset (1^n)\}, \quad j=\max\{k \ | \ \l_k\not\subset (1^n)\}.\]
Note that $1\leq i \leq j \leq l+1\leq n$. We associate to this walk the diagram $e_{i,j+2}$, which we note has height $p_j-2=j-1$.

Why is this a bijection? even a heuristic argument? 

Examples:

Consider the $J_{2,8}$-module $\Delta_{6,(6)}$, we give the following examples:
\begin{align*}
    0&-(1)-(2)-(3)-(4)-(3)-(4)-(5)-(6), &&\sim e_{4,5} &(ht&=-1)\\
        0&-(1)-(2)-(2,1)-(3,1)-(3)-(4)-(5)-(6), &&\sim e_{2,5} &(ht&=2).\end{align*}


\subsubsection{Many-cup, 1-D Specht Case}

Try two cup first? In the previous case, we made use of the fact that the sequence $p_k$ decreased exactly once. Here it will decrease more than once, will be hard to figure out the config of cups! 

\subsubsection{1-cup, non-1-D Specht Case}

In the 1D specht case, there is a unique ideal walk. Here, given a walk, we will need a way to determine the underlying ideal walk.

\subsubsection{Full Generality}
Hopefully can combine the previous two points

